% Исследовательский обзор (русский). Сборка: make research → thesis/build/research-ru.pdf
% Синтез reports/, research_notes/, docs/pedagogy/, docs/research/. Источники — SRC-NNN.
% Оформление — как у ВКР (Положение о ГИА Университета Иннополис, п. 6.8): 14 pt, чёрный Times, поля 25/20 мм.
\documentclass[a4paper,14pt]{extarticle}
\newcommand{\mainlang}{russian}
\newcommand{\docstyle}{thesis}
% Преамбула: шрифты (латиница, кириллица, арабский с огласовками), поля, библиография.
\usepackage{etoolbox}
\usepackage{fontspec}
\usepackage{polyglossia}
\setdefaultlanguage{\mainlang}
\ifdefstring{\mainlang}{english}{\setotherlanguages{russian,arabic}}{\setotherlanguages{english,arabic}}

% Times New Roman есть в macOS; на другой машине — Liberation Serif или Tempora.
\setmainfont{Times New Roman}
\newfontfamily\cyrillicfont{Times New Roman}
% Арабский — тот же шрифт, что рисует карточки бота (ADR-0024): огласовки и шадда не теряются.
\newfontfamily\arabicfont[Script=Arabic,Path=../src/muallim/presentation/telegram/assets/fonts/,
  Extension=.ttf,UprightFont=*-Regular,BoldFont=*-Bold,Scale=1.25]{ScheherazadeNew}
% \ar{…} — арабская вставка в строке; источник строки — в тексте рядом ([SRC-NNN]).
\newcommand{\ar}[1]{\textarabic{#1}}

\usepackage[left=3cm,right=1.5cm,top=2cm,bottom=2cm]{geometry} % [утвердить] по требованиям вуза
\usepackage{setspace}
\onehalfspacing
\usepackage{indentfirst}

\usepackage{graphicx}
\graphicspath{{../docs/architecture/diagrams/}{../docs/design/screens/}{figures/}}
\usepackage{booktabs,tabularx,longtable}
\usepackage{caption}
\usepackage{csquotes}
\usepackage[hidelinks]{hyperref}

% Библиография — тот же sources.bib, что у всего проекта (ключи SRC-NNN).
\ifdefstring{\mainlang}{english}
  {\usepackage[backend=biber,style=ieee,sorting=none]{biblatex}}
  {\usepackage[backend=biber,style=gost-numeric,sorting=none]{biblatex}}
\addbibresource{../docs/research/sources.bib}

% ID рисунков и таблиц из docs/thesis/figures.md: \label{fig:R-01}, \label{tab:T-01}.

% Данные для титульных листов всех документов. [утвердить] — с руководителем и учебным офисом.
\newcommand{\thesisProgramCode}{09.04.01}                     % [утвердить] код направления
\newcommand{\thesisProgramRu}{Информатика и вычислительная техника} % [утвердить]
\newcommand{\thesisProgramEn}{Computer Science}                % [утвердить]
\newcommand{\thesisSpecialtyRu}{Программная инженерия}
\newcommand{\thesisSpecialtyEn}{Software Engineering}
% Тема — предложение автора (2026-10-18; с ADR-0051 — с адаптацией модели), утверждает руководитель [утвердить].
\newcommand{\thesisTopicRu}{ИИ-тренажёр чтения по букварю «Муаллим ас-Сани»: адаптация открытой модели распознавания фонем к речи русскоязычных начинающих и оценка обратной связи}
\newcommand{\thesisTopicEn}{An AI Reading Tutor for the \emph{Muallim al-Sani} Primer: Adapting an Open Phoneme Recognition Model to Russian-Speaking Beginners and Evaluating Pronunciation Feedback}
\newcommand{\thesisShortRu}{ИИ-тренажёр «Муаллим ас-Сани» и адаптация модели}
\newcommand{\thesisShortEn}{AI Tutor for \emph{Muallim al-Sani} and Model Adaptation}
\newcommand{\thesisAuthorRu}{Исмаилов Абдирахим Ерланович}
\newcommand{\thesisAuthorEn}{Abdirakhim Ismailov}               % латиница — [утвердить] как в паспорте
\newcommand{\thesisAuthorEnFull}{Ismailov Abdirakhim Erlanovich}
\newcommand{\thesisSupervisorRu}{Аскарбекулы Нурсултан}
\newcommand{\thesisSupervisorEn}{Nursultan Askarbekuly}
\newcommand{\thesisUniversityRu}{Университет Иннополис}
\newcommand{\thesisUniversityEn}{Innopolis University}
\newcommand{\thesisCity}{Иннополис}
\newcommand{\thesisCityEn}{Innopolis}
\newcommand{\thesisYear}{2027}

\hypersetup{pdftitle={Научить машину слушать, как учитель букваря},pdfauthor={\thesisAuthorRu}}
\fancyhead[R]{\small\headfamily\color{textmuted}Исследовательский обзор · А.\,Е.~Исмаилов}
\setcounter{secnumdepth}{2}

\begin{document}
\thispagestyle{plain}
\begin{center}
{\headfamily\small\color{textmuted}\thesisUniversityRu{} · магистратура «\thesisSpecialtyRu» · исследовательский обзор к ВКР}\\[10pt]
{\headfamily\LARGE\bfseries\color{textaccent}Научить машину слушать, как учитель букваря\par}
\vspace{4pt}
{\headfamily\large\color{textaccent}Распознавание ошибок чтения, адаптация модели и оценка усвоения у начинающих читать Коран по «Муаллим ас-Сани»: состояние исследований и пробелы\par}
\vspace{10pt}
{\color{textrule}\rule{0.35\linewidth}{1pt}}\\[8pt]
{\large \thesisAuthorRu}\\[2pt]
{\small Научный руководитель: \thesisSupervisorRu}\\[2pt]
{\small\color{textmuted} Версия 1.0 · октябрь 2026}
\end{center}

\begin{keybox}[Аннотация]
Обзор сводит вместе семь областей знания, на которых стоит работа «\thesisTopicRu»: историю и
метод книги, фонетические трудности русскоязычных начинающих, исследования компьютерного обучения
произношению, состояние автоматической оценки чтения Корана, способы адаптировать речевую модель к новой
аудитории на малых данных, репетиторы, которые слушают чтение, вместе с моделями усвоения знаний, и
складывающуюся практику разработки ПО ИИ-агентами. Отдельно кратко описаны правовые и этические ограничения на сбор
голосовых данных. Доказательства поддерживают узкий тезис: автоматическая обратная связь может умножить
практику и первичную коррекцию для начинающих, но пока не может заменить учителя. Лучшая открытая
модель, Quran Muaalem, обучена на чтении мужчин-экспертов, и никто не измерял её на букварном чтении
тех, для кого арабский не родной, и никто не пробовал к ним её адаптировать, хотя методы для этого
--- дообучение низкого ранга, калибровка, отказ от ответа с гарантией --- существуют и работают на
малых данных, если оценку не испортить утечкой. Отслеживание знаний умеет оценивать усвоение навыка по
шумному выходу распознавателя, но ошибки распознавателя, измеренные против учителей, в такие модели ещё не
закладывали. Явная обратная связь на уровне отдельных звуков --- условие, при
котором обучение произношению работает лучше всего, и именно это нужно букварю. Разработка с ИИ-агентами
продуктивна, но склонна к подгонке под собственные тесты, поэтому нужны механические ограничения, а не
письменные правила. Обзор завершается картой пробелов и исследовательскими вопросами, которые из неё
следуют.

\smallskip
\textbf{Ключевые слова:} компьютерное обучение произношению; обнаружение и диагностика ошибок
произношения; чтение Корана; арабский как второй язык; адаптация модели; LoRA; калибровка;
конформное предсказание; отслеживание знаний; репетитор чтения; набор данных; design science; ИИ-агенты в
разработке ПО.
\end{keybox}

\section{Введение}

В конце XIX века ученик татарской школы учился читать по букварю и с учителем, который слышал каждый
слог. «Муаллим ас-Сани» («Второй учитель») Ахмад-Хади Максуди принадлежит этой традиции и используется
до сих пор: его переиздают в Татарстане и распространяет Духовное управление мусульман Казахстана
\cite{SRC-020,SRC-040}; сегодня это популярная в России и Казахстане методика обучения чтению на арабском языке. Взрослые, которые открывают его сегодня, часто говорят по-русски, начинают с
нуля и занимаются в одиночку между уроками. Им недостаёт слушающей половины метода.

Обзор отвечает на вопрос, может ли эту половину взять на себя машина и при каких условиях. Это
исследовательская основа магистерской работы, объект которой --- набор речевых данных и обученные на нём модели. Тренажёр, который
ведёт ученика по книге страница за страницей, служит инструментом сбора: с согласия учеников он собирает
букварное чтение, преподаватели его размечают, на этих данных открытая модель адаптируется к русскоязычным
начинающим, а поверх неё модель усвоения оценивает, какие навыки книги ученик уже освоил. Обзор построен
вокруг шести вопросов:

\begin{enumerate}[label=\textbf{В\arabic*.}]
  \item Что книга и традиция обучения по ней требуют от тренажёра?
  \item Что делает арабское чтение трудным для русско- и казахоязычных начинающих?
  \item Что известно об автоматической обратной связи по произношению и о том, какой она должна быть?
  \item Насколько хороши современные модели оценки чтения Корана и где они не проверены?
  \item Как приспособить такую модель к новой аудитории, когда данных мало, и как честно измерить выигрыш?
  \item Как по попыткам, которые слушает машина, оценить, что ученик уже усвоил?
\end{enumerate}

Седьмой, методический вопрос касается того, как создаётся сам тренажёр: что известно о разработке
исследовательского ПО ИИ-агентами и как удержать такую разработку честной.

\paragraph{Метод обзора} Обзор нарративный и обзорный (scoping), а не систематический. Поиск выполнен
в октябре 2026~г. по Google Scholar, arXiv, ACL Anthology, архиву ISCA, Crossref, сайтам издательств,
российским журналам из eLibrary, магазинам приложений и официальным страницам регуляторов. Каждый
источник получил постоянный ключ и оценку надёжности: A --- рецензируемые публикации и первоисточники,
B --- препринты, официальная документация и код авторов, C --- вторичные источники (блоги, страницы
магазинов), только для ориентира. Выходные данные ключевых научных источников сверены с Crossref или
arXiv. Полный реестр --- что взято из каждого источника и где это используется в коде --- хранится в
репозитории проекта \cite{SRC-038}.

\section{Букварь и его метод}

\subsection{Две книги Максуди}

Максуди (1868--1941) написал два букваря, которые в интернете постоянно путают. «Мөгаллим әүвәл»
(«Первый учитель», 1892) учил татарской грамоте на арабской графике; звуко-слоговой метод и более
тридцати изданий к 1918~г. сделали его одной из самых издаваемых татарских книг своего времени
\cite{SRC-015,SRC-019}. «Муаллим ас-Сани» продолжает его: от знания букв к правильному чтению Корана с
основами \emph{таджвида}. Книгу обычно относят к середине 1890-х, но год первого издания по первичному
библиографическому источнику не подтверждён \cite{SRC-018}. Современные тексты восходят к казанскому
изданию 1917~г. \cite{SRC-024}. Издания различаются объёмом и содержанием, поэтому в работе
зафиксировано одно конкретное --- альметьевское издание 2010~г., вышедшее при Духовном управлении
мусульман Татарстана \cite{SRC-040}, --- и каждое упражнение привязано к его странице и строке.

\subsection{От названий букв к звукам}

Татарская школа знала три способа обучения грамоте \cite{SRC-017}. Старый буквослагательный метод
(\emph{усул-и хәрфия}) заставлял сначала заучивать названия букв --- \emph{алиф, ба, та} --- и только
потом складывать слоги. Слоговой метод (\emph{усул-и мәддия}) сразу учил слогам. Звуковой метод
джадидов (\emph{усул-и саутия}) учил букву как звук и осознанно собирал из звуков слоги и слова.
Переходу к звуковому методу приписывают сокращение обучения чтению «с нескольких лет до одного
учебного года» \cite{SRC-016}. Дамиров добавляет важную для мусульманской аудитории деталь:
реформатор З.~Расулев обосновывал звуковой метод традицией самого \emph{таджвида}, а не европейской
педагогикой \cite{SRC-016}.

Новаторство Максуди, по Фахрутдиновой и Ибрагимову, --- в трёх решениях: отказе от заучивания
названий букв, внимании к артикуляции и звучанию, постепенности и осмысленности \cite{SRC-015}. Его
пятнадцать дидактических принципов, опубликованных в 1906~г., сводятся к четырём группам. В
таблице~\ref{tab:principles-ru} они сопоставлены с решениями для тренажёра.

\begin{table}[ht]
\centering\small
\caption{Дидактические принципы Максуди \cite{SRC-015} и что из них следует для цифрового тренажёра.}
\label{tab:principles-ru}
\begin{tabularx}{\linewidth}{@{}P{4.4cm}L@{}}
\toprule
Принцип & Что следует для тренажёра \\
\midrule
Постепенность и преемственность & Уроки в порядке книги; строка считается освоенной, только когда
прочитана верно; маленькие шаги (урок не длиннее дюжины строк). \\
Индивидуализация и активность & Свой темп и своя история попыток; повтор слабых строк; сколько угодно
попыток. \\
Наглядность и апперцепция & Эталонное аудио до чтения; буква урока выделена, как в книге; подсказки,
связывающие новый звук со знакомыми русскими звуками. \\
Практика и непрерывный контроль & Каждая попытка проверяется сразу; каждая ошибка объясняется; можно
сразу повторить. \\
Гуманность, без наказаний & Без штрафов и рейтингов; серия дней --- поощрение, а не угроза. \\
Звук, а не название буквы & Карточка учит звук с огласовкой; название --- только справка. \\
\bottomrule
\end{tabularx}
\end{table}

\subsection{Порядок букв и цикл урока}

Два независимых источника и издание 2010~г. сходятся: книга не следует алфавиту. Буквы вводятся
группами, начиная с хамзы и \ar{ر ز م} --- букв простой, в основном несоединяемой формы, из которых сразу
складываются слоги вроде \ar{زَرَ} и \ar{مَرَ} \cite{SRC-021,SRC-040}. Дальше идут \ar{ت ن},
\ar{ي ب}, \ar{ك ل} и так до \ar{ظ}; затем долгие гласные, \emph{шадда} и \emph{танвин}, артикль и
\emph{сукун}, правила \emph{мим} и \emph{нун}, особенности письма Корана, остановки и, наконец, короткие
суры. Тот же порядок встречается в узбекском издании и в современном курсе научного руководителя
\cite{SRC-070}, что говорит о том, что порядок принадлежит самому Максуди.

Все найденные курсы следуют одному циклу --- цифровой форме \emph{таляккы}: услышать образец,
повторить (в видеокурсе каждую строку трижды), прочитать самостоятельно, получить исправление и идти
дальше только после безошибочного чтения \cite{SRC-021,SRC-022,SRC-058}. Два наблюдения практиков
определяют дизайн. Преподаватели не советуют учиться по транскрипции: «на бумаге невозможно передать
звук» \cite{SRC-021,SRC-022}. А казахское издание прямо оговаривает, что книга не заменяет живого
учителя. Значит, тренажёр должен ставить на первое место звук эталона, держать транскрипцию как скрытую
подсказку и представлять себя партнёром по практике, а не заменой учителя.

\section{Ученики}

Целевые ученики --- русскоязычные взрослые в России и Казахстане, которые не знают арабских букв и
звуков, отсутствующих в русском. Рецензируемых исследований их ошибок чтения не найдено. Таблица
\ref{tab:difficult-ru} объединяет методические источники для русскоязычных с исследованиями носителей
других языков \cite{SRC-023,SRC-034,SRC-035}. Фарингальные, эмфатические и межзубные согласные трудны
почти для всех; долгота гласных и удвоение трудны потому, что в русском они не различают смысл.

\begin{table}[ht]
\centering\small
\caption{Типичные трудности русскоязычных начинающих (методические источники и исследования других L1).}
\label{tab:difficult-ru}
\begin{tabularx}{\linewidth}{@{}P{3.2cm}L@{}}
\toprule
Звуки & Что обычно делают ученики \\
\midrule
\ar{ح}, \ar{ع} & \ar{ح} превращается в русское [х] или сливается с \ar{خ}, \ar{ه}; \ar{ع} --- в гортанную
  смычку или пропадает в гласном. \\
\ar{ق}, \ar{غ}, \ar{خ} & \ar{ق} произносят «передним» [к]; \ar{غ} и \ar{خ} путают. \\
\ar{ص}, \ar{ض}, \ar{ط}, \ar{ظ} & Эмфатические теряют твёрдость и звучат как [с д т з]. \\
\ar{ث}, \ar{ذ}, \ar{ظ} & Межзубные становятся [с]/[т], [з]/[д]. \\
Долгие гласные, \emph{шадда} & Долгие укорачивают; удвоенные согласные не удваивают. \\
\ar{ر} & Смягчается перед гласными переднего ряда. \\
\bottomrule
\end{tabularx}
\end{table}

Букварь эти путаницы предвидит. Его строки «обратите внимание» противопоставляют \ar{ق}/\ar{ك},
\ar{ث}/\ar{س}, \ar{ص}/\ar{س}, \ar{ط}/\ar{ت}, \ar{ح}/\ar{خ}, \ar{ع}/\ar{غ}, \ar{ض}/\ar{د} и
\ar{ذ}/\ar{ز} --- в современных терминах это минимальные пары \cite{SRC-040}. Для казахоязычных из
фонологии следует проверяемая гипотеза: в казахском есть увулярные [q] и [\ipa{ʁ}], поэтому \ar{ق} и
\ar{غ} им, возможно, даются легче, чем русским, а остальные трудности общие. Пока нет данных, это
только гипотеза.

\section{Компьютерное обучение произношению}

\subsection{Помогает ли автоматическая обратная связь?}

Самые сильные доказательства дают метаанализы в обучении второму языку. Нго, Чен и Лай объединили
исследования распознавания речи в обучении английскому произношению и получили средний общий эффект
$g = 0{,}69$; эффект больше для отдельных звуков, чем для просодии, и больше при явной корректирующей
обратной связи, чем при неявной (просто транскрипт) \cite{SRC-025}. Обучение произношению вообще даёт
большой эффект \cite{SRC-105}, так что автоматическая обратная связь не лучше обучения --- это способ
дать его больше. Систематический обзор тридцати исследований Амрате и Цай добавляет два
предостережения: распознаватели часто ошибаются на речи с акцентом и дают ложные замечания на понятное
произношение, поэтому рекомендуется использование под контролем учителя или партнёра, способного
поправить систему \cite{SRC-026}. Типичное исследование квазиэкспериментальное, с несколькими
десятками участников в течение нескольких недель --- реалистичная планка для магистерского пилота.

Две соседние линии исследований поддерживают функции за пределами основного цикла. Тренировка
восприятия на голосах разных дикторов (HVPT) даёт средний или большой и устойчивый эффект на
различение звуков \cite{SRC-044} и меньший перенос на произношение \cite{SRC-045}. Распределённая
практика --- один из самых надёжных результатов науки об обучении \cite{SRC-107}; это обосновывает
повторение слабых строк по расписанию и отложенный тест.

\subsection{Какой должна быть обратная связь?}

Исследования корректирующей обратной связи на уроках показали, что простое повторение правильной формы
(рекаст) редко приводит к самоисправлению, а подсказки, указывающие на ошибку и просящие попробовать
снова, --- приводят \cite{SRC-053}. Работы о компьютерном обучении произношению советуют разбирать
немногие важные ошибки, а не все подряд \cite{SRC-054}. Обзоры формирующей обратной связи требуют
конкретности, простоты и небольших порций, а подробности --- по запросу \cite{SRC-055}, и называют
самым полезным элементом ответ на вопрос «что делать дальше», тогда как похвала личности почти ничего не
даёт \cite{SRC-056}. Коранические методики Икро и \emph{таляккы} приходят к тому же из практики:
учитель показывает, где ошибка, а ученик исправляет её, перечитывая \cite{SRC-057,SRC-058}.
Коммерческие приложения подсвечивают ошибки прямо в тексте \cite{SRC-046}.

Вместе эти источники подсказывают единицу обратной связи из трёх частей: \emph{что} не так и где (слог
отмечен в строке), \emph{как} сказать (подсказка по артикуляции на языке ученика) и \emph{что дальше}
(послушай, скажи слог медленно, прочитай строку снова). Они же подсказывают показывать за раз
немного ошибок и ставить явные ошибки, например неверную букву, раньше тонких, например долготы
гласного \cite{SRC-052}.

\subsection{Telegram как среда}

Telegram --- повседневный мессенджер целевой аудитории; голосовые сообщения в нём работают без
установки отдельного приложения. Четырёхнедельное исследование показало, что Telegram-канал улучшил
английское произношение иранских студентов \cite{SRC-092}; контролируемых исследований Telegram-бота с
автоматической проверкой произношения не найдено. У платформы свои правила: данные для машинного
обучения можно собирать только с явного, активного и отзываемого согласия пользователя, а политика
конфиденциальности обязательна \cite{SRC-033}.

\section{Автоматическая оценка чтения Корана}

\subsection{Задача и метрики}

Обнаружение и диагностика ошибок произношения сравнивают то, что сказал ученик, с тем, что требует
текст, на уровне фонем. Стандартная оценка, принятая в L2-ARCTIC и в коранических бенчмарках, относит
каждую фонему эталона к верному принятию, верному отклонению (настоящая ошибка, которую система
нашла), ложному принятию (пропущенная ошибка) или ложному отклонению (ложная тревога). Precision ---
это TR/(TR+FR), recall --- TR/(TR+FA), F1 --- их гармоническое среднее; точность диагноза показывает,
назвала ли система правильную замену \cite{SRC-063,SRC-064}. Для начинающего эти два вида ошибки не
равноценны. Пропущенная ошибка замедляет обучение, а ложная тревога на верном чтении подрывает доверие
и мотивацию.

\subsection{Бенчмарки}

QuranMB, представленный на Interspeech 2025, стал первым открытым бенчмарком обнаружения ошибок
произношения в арабском на материале Корана \cite{SRC-004}. Его базовые системы имели высокую полноту и
очень низкую точность. В соревновании Iqra'Eval системы участников достигли F1 около 0,40--0,44
\cite{SRC-043}. В конкурсе IQRA 2026 на QuranMB.v2 лучшая система получила F1 0,72 по фонемам против
0,44 у базовой \cite{SRC-031,SRC-086}. Побеждают распознаватели фонем на самообучаемых речевых
энкодерах с CTC --- то же семейство, что и Muaalem. Организаторы прямо пишут, что в данных нет
изучающих арабский как второй язык, и призывают собирать такие данные \cite{SRC-031}.

\subsection{Quran Muaalem}

Quran Muaalem --- самая полная открытая система диагностики ошибок чтения Корана \cite{SRC-002,SRC-003}.
Это дообученный энкодер Wav2Vec2-BERT с одиннадцатью CTC-головами: одна --- для фонем в специальной
фонетической записи Quran Phonetic Script, десять --- для артикуляционных признаков (\emph{сифат}):
звонкости, твёрдости, \emph{калькаля} и др. Обучающий корпус --- 848 часов и около 286 тыс. фрагментов
от 22 чтецов-экспертов, размеченных в основном автоматически. Заявленные результаты --- в
таблице~\ref{tab:muaalem-ru}.

\begin{table}[ht]
\centering\small
\caption{Заявленное качество Quran Muaalem \cite{SRC-002}.}
\label{tab:muaalem-ru}
\begin{tabularx}{\linewidth}{@{}Lr@{}}
\toprule
Мера & Значение \\
\midrule
Доля ошибок по фонемам (PER), тест авторов (чтение экспертов) & 0,54\,\% \\
Средняя доля ошибок по всем одиннадцати уровням, тот же тест & 0,21\,\% \\
PER по фонемам, qdat\_bench (159 записей с настоящими ошибками) & 5,78\,\% \\
Средний F1 по таджвиду на qdat\_bench (три правила) & 0,758 \\
F1 для \emph{ихфа} на qdat\_bench & 0,453 \\
\bottomrule
\end{tabularx}
\end{table}

Десятикратный разрыв между тестом на чтении экспертов и записями с настоящими ошибками --- главное число
для тренажёра начинающих. Ограничения называют сами авторы: обучающие данные --- только мужчины-эксперты
и только риваят Хафс, а qdat\_bench размечен одним человеком на одном аяте \cite{SRC-002}. Фонетизатор
рассчитан на текст Корана, поэтому для отдельных букв и слогов букваря нужны канонические строки фонем,
а медленное послоговое чтение лежит вне распределения обучающих данных. Пакет также отдаёт уверенность
по каждой фонеме и вероятности признаков \cite{SRC-062}, что открывает путь к исходу «не уверен» вместо
ложного вердикта.

\subsection{Данные об учениках}

Данных о начинающих-неарабах мало. Ближайший открытый ресурс --- краудсорсинговый набор примерно из семи
тысяч чтений неарабов более чем из одиннадцати стран --- создан в Университете Иннополис \cite{SRC-032};
носителей русского и казахского в нём нет. Работы той же группы о проверке заучивания в Telegram-боте
\cite{SRC-067} и о таксономии событий чтения \cite{SRC-068} --- ближайшие соседи этой работы, с важным
различием. Там проверяется, \emph{что} прочитано: какие аяты, в каком порядке, с пропусками и
вставками, по транскрипту без огласовок; огласовки и \emph{таджвид} прямо вынесены за рамки. Тренажёр
букваря должен проверять, \emph{как} произнесён слог. Таксономия всё равно полезна: она отделяет
неисправленные ошибки от самоисправлений и безобидных событий, а на записях начинающих это различие
понадобится.

\subsection{Приложения}

Tarteel, лидер рынка, ориентирован на заучивание и отмечает ошибки на уровне слов; он подсвечивает их в
тексте --- приём, который стоит взять \cite{SRC-046}. Новая волна приложений обещает ИИ-проверку на
уровне букваря, но эти обещания взяты из описаний в магазинах и независимо не проверены. Цифровые версии
«Муаллим ас-Сани» существуют в виде приложения и открытой веб-платформы \cite{SRC-005,SRC-006}, а курс
научного руководителя идёт в том же порядке букв, но тренирует слух, а не произношение \cite{SRC-070}.
Инструмента на русском, казахском или узбекском, который проверял бы чтение букварного уровня голосом,
не найдено.

\section{Адаптация модели к новой аудитории}

\subsection{Дообучение на речи изучающих язык}

Модели обнаружения ошибок на самообучаемых речевых энкодерах обычно получают дообучением с CTC на
размеченной речи изучающих язык. Сюй и соавторы показали на L2-ARCTIC, что дообученный wav2vec~2.0 лучше
прежних систем и сохраняет качество, когда размеченных данных мало \cite{SRC-116}; Muaalem построена так
же, только дообучалась на экспертах \cite{SRC-002}. Для нескольких часов речи полное дообучение
600-миллионной модели рискованно: оно переобучается и может испортить то, что модель уже умеет. Метод
LoRA замораживает исходные веса и обучает малые матрицы низкого ранга, добавленные к слоям внимания;
обучаемых параметров на порядки меньше, а исходную модель можно вернуть, просто отключив добавки
\cite{SRC-111}. Дообучение на новой задаче может стирать старые умения --- катастрофическое забывание
\cite{SRC-118}, --- поэтому адаптированную модель нужно проверять и на прежнем домене, здесь --- на чтении
экспертов.

\subsection{Где взять ошибки}

У начинающих ошибок больше, чем у экспертов, но для обучения их всё равно мало, а отдельные пары звуков
встречаются редко. В литературе по MDD данные пополняют синтетическими ошибками: в эталонной
транскрипции заменяют фонемы, как будто ученик ошибся \cite{SRC-119}; стандартный приём распознавания
речи --- изменение скорости записи \cite{SRC-120}. Букварь даёт более естественный источник: его
минимальные пары (\ar{ق}/\ar{ك}, \ar{ص}/\ar{س}, \ar{ح}/\ar{خ} и др.) означают, что верное чтение одной
карточки --- настоящая запись настоящей ошибки для парной. Такое пополнение использует аудио людей, а не
синтез: синтетический арабский голос модель распознаёт плохо (предварительная проверка проекта).

\subsection{Уверенность, калибровка и «не уверен»}

Второй путь --- не менять модель, а научиться ей не доверять там, где она ошибается. Классическая мера
Goodness of Pronunciation превращает апостериорные вероятности фонем в оценку произношения
\cite{SRC-115}; поверх таких признаков можно обучить небольшой верификатор на сотнях размеченных примеров.
Нейросети, однако, переуверены, и их вероятности нужно калибровать; стандартная мера --- ожидаемая ошибка
калибровки \cite{SRC-112}. Выборочная классификация разрешает системе отказаться от ответа и описывает
компромисс кривой «риск --- покрытие» \cite{SRC-117}. Конформное предсказание даёт гарантию без допущений
о распределении при перестановочности данных \cite{SRC-113}, а конформный контроль риска позволяет
выбрать порог так, чтобы ожидаемая доля ложных отклонений не превышала заданный уровень \cite{SRC-114}.
Для тренажёра это прямой ответ на главную опасность, о которой предупреждают обзоры обучения
произношению: ложное «не засчитано» на верном чтении.

\subsection{Как не обмануть себя}

С малыми данными легко получить завышенный результат. Подбор гиперпараметров на тех же данных, на которых
сообщается качество, смещает оценку вверх; нужна вложенная кросс-валидация или отдельный тест
\cite{SRC-121}. Капур и Нараянан нашли утечку данных в сотнях ML-исследований и предложили её таксономию
\cite{SRC-122}; для речи главная утечка --- один и тот же диктор в обучении и проверке. Отсюда правила
работы: разбиение по ученикам, тест, замороженный до обучения, разметка теста двумя учителями и один прогон на нём.
Те же требования группа научного руководителя сформулировала для исследований, где модели подбирают
ИИ-агенты \cite{SRC-067}. Набор, который будет опубликован, сопровождается карточкой по Гебру и
соавторам \cite{SRC-123}.

\section{Репетиторы, которые слушают чтение, и измерение усвоения}

\subsection{Освоение чтения, а не только произношение}

«Муаллим ас-Сани» --- букварь: ученик учится видеть букву с огласовкой и превращать её в звук. Ближайшие к
работе системы поэтому не тренажёры произношения иностранного языка, а репетиторы чтения, которые слушают,
как ребёнок читает вслух. Project LISTEN в университете Карнеги-Меллон двадцать пять лет развивал
репетитор, который по распознаванию речи решал, верно ли прочитано слово, помогал на трудных словах и
следил за прогрессом \cite{SRC-125}. Нидерландский проект DART построил репетитор для первоклассников и
разделил два вопроса: насколько распознавание согласуется с людьми ($\kappa$, MCC) и как учатся дети ---
эксперименты с обратной связью на 525 и 752 детях и прогресс чтения по логам \cite{SRC-130,SRC-131}.
Французский Lalilo адаптирует самообучаемые модели к фонемам начинающих читателей и симулирует ошибки чтения,
чтобы модели было на чём учиться; этому посвящена докторская диссертация Желен \cite{SRC-132,SRC-133}.
Все три системы работают с детьми, которые читают на родном языке; взрослых, которые учатся читать чужое
письмо, среди них нет.

Ошибки начинающего читателя бывают двух родов. Он может прочитать не тот знак --- другую огласовку, не ту из
похожих на вид букв, --- и может, прочитав знак верно, неверно произнести звук. Исследования ранних читателей
арабского находят больше ошибок в гласных, чем в согласных, и путаницу букв, различающихся точками, но это
дети-носители. Для взрослых неарабов такого разделения нет ни в одном наборе данных.

\subsection{Отслеживание знаний}

Байесовское отслеживание знаний (BKT) моделирует каждый навык как скрытое состояние «усвоен / не усвоен» и
обновляет вероятность усвоения после каждой попытки; параметры угадывания и оплошности допускают, что
усвоивший навык иногда ошибается, а не усвоивший --- угадывает \cite{SRC-124}. Обзор Пеланека сравнивает BKT с
логистическими моделями и рекомендует оценивать модели ученика по точности предсказания следующей попытки
(AUC) и калибровке \cite{SRC-126}. Нейросетевое отслеживание знаний \cite{SRC-127} точнее на больших данных,
но на десятках учеников оправдано только как разведка. Модель аддитивных факторов строит кривые обучения по
навыкам и проверяет, удачно ли выбраны сами навыки \cite{SRC-128}; регрессия периода полураспада памяти
моделирует забывание по логам \cite{SRC-129}.

Ближайшая к работе публикация --- Бек и Сисон \cite{SRC-125}: они оценивали усвоение 80 соответствий
«буква--звук» у детей по выходу распознавателя речи, который сам ошибается. Шум распознавателя там поглощался
параметрами угадывания и оплошности; против учителей он не измерялся. Если же ложные принятия и отклонения
распознавателя по каждой фонеме измерены на размеченных данных, их можно заложить в модель явно, а адаптация
распознавателя должна делать оценку усвоения точнее. Этого не проверяли.

\subsection{Оценка по усвоению}

Систематический обзор образовательных рекомендательных систем показал, что их чаще оценивают по кликам и
оценкам, а по результатам обучения --- лишь примерно в трети работ \cite{SRC-135}. Рекомендательная система
OBER хранит результаты обучения в схеме данных и сравнивает алгоритмы по усвоению; в рандомизированном
испытании на 5\,700 учениках алгоритм, дольше удерживавший учеников, не дал лучшего усвоения
\cite{SRC-134}. Для тренажёра это значит: судить нужно по тому, что ученик усвоил, а не по тому, сколько он
занимался.

\begin{keybox}[Вывод]
Репетиторы, которые слушают чтение, существуют, но для детей и родного письма. Отслеживание знаний умеет
работать с шумным распознаванием, но не учитывает ошибки распознавателя, измеренные против учителей. Уроки
«Муаллим ас-Сани» дают готовую карту навыков. Отсюда дополнительный вопрос работы: как инструмент влияет на
динамику освоения материала (RQ-3); измерять её позволяет оценка усвоения по навыкам книги с учётом ошибок
распознавателя.
\end{keybox}

\section[Разработка исследовательского ПО с ИИ-агентами]{Разработка исследовательского ПО\newline с~ИИ-агентами}

\subsection{От цикла агента к обвязке}

ИИ-агент работает циклом: рассуждает, действует инструментом, наблюдает результат \cite{SRC-108}.
SWE-agent показал, что важен интерфейс между агентом и компьютером: линтер, отклоняющий синтаксически
неверные правки в момент правки, заметно увеличил долю решённых задач \cite{SRC-109}. Практики теперь
называют всё вокруг цикла \emph{обвязкой} (harness): инструкции, хуки, субагенты, тесты, линтеры, журналы.
Бёкелер различает \emph{направляющие} (guides), которые задают курс до действия, и \emph{датчики}
(sensors), которые наблюдают и исправляют после, а также вычислительные проверки вроде проверки типов и
инференциальные вроде ревьюера на основе LLM \cite{SRC-029}. Инженерные заметки Anthropic советуют
внешнее состояние, одну функцию за сессию, сквозную проверку и явный запрет править тесты
\cite{SRC-028}. Рецензируемого определения «harness engineering» пока нет; разведочное исследование
2\,853 репозиториев показало, что преобладают файлы контекста, а скиллы и субагенты редки
\cite{SRC-073}.

\subsection{Что говорят данные}

Прямое отношение к работе имеют три результата. Первый: агенты подгоняют оценку. В ImpossibleBench, где
спецификация противоречит тестам и любое прохождение означает обман, передовые модели часто всё равно
«проходили», причём у моделей Claude большая часть обмана сводилась к правке тестов; тесты только для
чтения закрывали этот путь \cite{SRC-027}. Подгонка под награду, выученная при обучении, может даже
обобщаться в саботаж \cite{SRC-082}. Второй: письменные правила слабы. Файлы контекста вроде
\texttt{AGENTS.md} не повышали, а иногда снижали успешность при росте стоимости \cite{SRC-074}; в
другом исследовании они сэкономили время и токены при той же успешности \cite{SRC-075}; абляция на 288
прогонах не нашла изменения корректности \cite{SRC-077}; факторное исследование 1\,650 сессий показало,
что соблюдение инструкций падает внутри сессии \cite{SRC-076}. Третий: механическая обратная связь
работает. Замечания статического анализа, возвращённые модели, заметно снизили число проблем
безопасности за десять итераций \cite{SRC-078}; агенты способны писать полезные тесты свойств
\cite{SRC-080}; генерация тестов на основе мутантов в Meta дала тесты, которые инженеры приняли
\cite{SRC-079}.

Картину дополняют два предостережения. Опытные разработчики в рандомизированном исследовании работали
с ИИ на 19\,\% медленнее, считая, что ускорились \cite{SRC-030}, так что продуктивность нужно мерить по
журналам, а не по ощущениям. А симулированные ученики на основе языковых моделей систематически
расходятся с настоящими \cite{SRC-083,SRC-085}: агент, «играющий ученика», может найти проблемы
интерфейса, но не заменяет пилот.

\begin{keybox}[Вывод]
Важные правила должны стать проверками, которые обеспечивает среда: замороженные тесты в фазе
реализации, архитектурные контракты, проверка типов, защищённые зоны человека и журнал каждой
блокировки. Этот принцип работа применяет к собственной разработке и к защите своих экспериментов; он
описан в методике как средство достоверности.
\end{keybox}

\section{Правовые и этические ограничения}

Голосовой тренажёр чтения Корана обрабатывает данные персональные, потенциально биометрические и
указывающие на религиозные убеждения. В России разработчик публичного бота --- оператор персональных
данных по 152-ФЗ; с сентября 2025~г. согласие оформляется отдельным документом, а база с данными
граждан должна находиться в России \cite{SRC-036}. Казахстан требует хранения данных на своей
территории и согласия на автоматизированную обработку \cite{SRC-037}. Правила Telegram запрещают
собирать данные для машинного обучения без явного, активного и отзываемого согласия \cite{SRC-033}.
Поэтому защитимый дизайн отделяет согласие на работу сервиса от согласия на запись голоса, никогда не
хранит аудио без согласия, не строит голосовых отпечатков, использует псевдонимы во всех
исследовательских материалах, допускает к исследованию только взрослых и позволяет ученику удалить всё.
Тренажёр должен использовать только арабский текст мусхафа и собственные нейтральные пояснения;
русский текст книги защищён авторским правом и в продукт не попадает. Вопросы оператора данных и
этической экспертизы --- институциональные и решаются университетом.

\section{Синтез: пробелы и исследовательские вопросы}

Таблица~\ref{tab:gaps-ru} показывает, что известно, чего не хватает и как отвечает работа.

\begin{table}[ht]
\centering\small
\caption{Карта доказательств и пробелов.}
\label{tab:gaps-ru}
\begin{tabularx}{\linewidth}{@{}P{2.6cm}LLP{2.6cm}@{}}
\toprule
Область & Известно & Не хватает & В работе \\
\midrule
Букварь & Звуковой метод, принципы, порядок букв \cite{SRC-015,SRC-016,SRC-017} & Изучение «Муаллим
ас-Сани» как букваря Корана; цифровая адаптация с проверкой & Артефакт A \\
Данные & Беглое чтение аятов экспертами \cite{SRC-004,SRC-043}; неарабы без русских \cite{SRC-032} &
Букварное чтение начинающих, изучающих арабский как иностранный; ошибки чтения отдельно от ошибок произношения & RQ-1, артефакт A \\
Модели & Сильны на экспертах; F1 0,72 на бенчмарке \cite{SRC-002,SRC-031} & Точность на начинающих,
доля ложных отклонений, согласие с учителями & RQ-1 \\
Адаптация & LoRA, калибровка, конформные гарантии \cite{SRC-111,SRC-112,SRC-114} & Система из готовых
моделей и компонентов, собранная agentic loop-engineering, с честной оценкой на новых учениках против базовых
подходов & RQ-2 \\
Усвоение & BKT и шумное распознавание у детей \cite{SRC-124,SRC-125} & Учёт ошибок распознавателя,
измеренных против учителей; карта навыков букваря & инструмент RQ-3 \\
Ученики & Трудные звуки у носителей других языков \cite{SRC-034,SRC-035} & Как русскоязычные взрослые
осваивают навыки книги с инструментом, что забывают & RQ-3 (доп.) \\
Обратная связь & Средний эффект; лучше всего явная и на звуках \cite{SRC-025,SRC-026} & Эффект для
арабских букв; испытания букварных методик & разведка, PhD \\
Агенты & Подгонка тестов; правила забываются; проверки помогают \cite{SRC-027,SRC-076,SRC-078} &
Данные о целом проекте с обвязкой & методика \\
\bottomrule
\end{tabularx}
\end{table}

Из этой карты вытекают гипотеза и вопросы работы. Гипотеза: собрав небольшой курированный речевой датасет, можно оценить существующие решения и с помощью agentic loop-engineering разработать эффективный инструмент для внеклассной практики начинающих. RQ-1: насколько надёжно существующие модели распознавания и анализа речи обнаруживают ошибки чтения и произношения у начинающих изучать арабский язык и чтение Корана? RQ-2: насколько применение agentic loop-engineering для выбора и комбинирования существующих моделей и компонентов повышает точность обнаружения ошибок чтения и произношения по сравнению с базовыми подходами? RQ-3 (дополнительный): в какой степени использование разработанного инструмента во внеклассной практике влияет на динамику освоения учебного материала начинающими, обучающимися с преподавателем? Эффект объяснений
остаётся разведкой и вопросом докторантуры, а обвязка агента --- средством достоверности.

\section{Заключение}

Литература оставляет для этой работы ясное и скромное место. Метод «Муаллим ас-Сани» совместим с
автоматической обратной связью: он звуковой, пошаговый и построен на немедленном исправлении. Условия,
при которых обучение произношению работает лучше всего, --- отдельные звуки и явная обратная связь,
--- это и есть условия букваря. Модель существует и открыта. Не хватает измерения: как часто модель
права насчёт русскоязычного взрослого, впервые читающего слог, и как часто она отклоняет верное чтение;
насколько точнее можно сделать проверку, собирая её с помощью agentic loop-engineering из готовых моделей на
собственных данных новой аудитории; и как меняется освоение материала с таким инструментом. Небольшой честно размеченный набор данных, замороженный тест и прозрачная таблица precision и
recall относительно учителей до и после agentic loop-engineering стоят больше, чем ещё одно заявление об
«ИИ-проверке». Этот же датасет, тест и карта навыков становятся отправной точкой докторской работы о том, как
учить ИИ и учиться с ИИ. В том же духе и
методику разработки нужно судить по тому, что показывают журналы, а не по тому, насколько быстрой она
казалась.

\printbibliography[heading=bibintoc,title={Список источников}]
\end{document}
